\section{Introduction}\label{sec:introduction}
For a real vector $x=(x_{1-N},\dots,x_{N-1})$ let $T_N(x)=(x_{i-j})_{0\le i,j<N}$
be the real Toeplitz matrix of order $N$ with symbol $x$. Throughout, norms
and inner products of matrices are Frobenius norms and inner products. Put
\begin{equation}\label{eq:FN}
 F_N(x,y)=2\|T_N(x)\|_F^2\|T_N(y)\|_F^2
 -2\langle T_N(x),T_N(y)\rangle_F^2
 -\|[T_N(x),T_N(y)]\|_F^2,
\end{equation}
a homogeneous quartic in $4N-2$ real variables. The B\"ottcher--Wenzel
inequality implies that $F_N$ is nonnegative. L\'aszl\'o conjectured that
this nonnegativity can always be expressed as a sum of squares of real
quadratic forms \cite[Conjecture~15]{laszlo2012}; the question is problem
MI-15 in the OpenProblemsInNLA catalogue \cite{catalogue2026}. The
following theorem answers it negatively.

\begin{theorem}\label{thm:negative}
For every integer $N\ge\ObstructionOrder$, the polynomial $F_N$ is not a finite
sum of squares of real homogeneous quadratic forms. Nevertheless,
$F_N(x,y)\ge0$ for all real symbol vectors $x,y$ and every $N\ge1$.
\end{theorem}

A \emph{sum of squares} (SOS) always means a finite sum of squares of
real homogeneous quadratic forms, with arbitrary real coefficients and
any number of terms. For a quartic form, requiring the squares to be
homogeneous quadratic forms loses nothing: if a form of degree four
equals $\sum_jp_j^2$ for real polynomials $p_j$, then neither the squares
of the highest-degree homogeneous parts of the $p_j$ nor those of their
lowest-degree parts can cancel, so every $p_j$ is a quadratic form. The
nonnegativity assertion is classical (Section~\ref{sec:context}). The
threshold is sufficient, not optimized; the first order at which a sum of
squares fails is unknown.

\subsection{The obstruction in outline}
Let $\calB_m$ be the restriction of $F_{2m}$ to the $m$ outermost nonzero
diagonals on each side, all other symbol variables being zero. After
relabelling by depth, the same polynomial is a restriction of $F_N$ for
every $N\ge2m$ (Lemma~\ref{lem:corner}). Substitution preserves a sum of
squares, so a single obstruction for $\calB_m$ transfers to every larger
order.

\emph{What is tested.} By Lemma~\ref{lem:frame}, $\calB_m$ is a sum of
squares exactly when $\calB_m=z^TQz$ for a positive semidefinite Gram
matrix $Q$ on the wedge coordinates $z_{ab}=x_ay_b-x_by_a$; such a $Q$
is determined only up to the Pl\"ucker relations among the wedges. The
test is a linear functional $\Phi_m$ of $Q$ that is nonnegative whenever
$Q\succeq0$. To form it, average $Q$ over two symmetries of the corner,
which keeps it positive semidefinite, and split the result into three
positive semidefinite blocks. Substituting monomials in two complex
variables for the wedge coordinates turns each block into a generating
kernel in four complex variables. A fixed combination of these kernels,
evaluated at all pairs of $1024$ points of $\C^2$, is a positive
semidefinite Hermitian $1024\times1024$ matrix, and $\Phi_m$ is its
quadratic pairing with a fixed real vector $c$.

\emph{At which points.} The points and the coefficients are fixed once
and for all:
\begin{equation}\label{eq:witness}
\begin{gathered}
 I=\{(r,j,\sigma):r\in\{1,2\},\ 1\le j\le256,\ \sigma\in\{-1,1\}\},\\
 a_\ell=r+i\sigma j,\quad b_\ell=r-i\sigma j,\quad
 c_\ell=\sigma\hat c_j,\qquad
 \hat c_j=\left\lfloor\frac{2^{43}j}{(128^2+j^2)^2}\right\rfloor.
\end{gathered}
\end{equation}
Thus $|I|=1024$ and $\|c\|_1=859652168<2^{30}$. For a dyadic scale
$\epsilon=2^{-k}$ and a corner of depth $m=\epsilon^{-2}$, the point
attached to $\ell=(r,j,\sigma)$ is
$(ie^{-\epsilon a_\ell},-ie^{-\epsilon b_\ell})$. It weights a diagonal at
depth $p$ by $e^{-\epsilon pr}$ and therefore probes depths of order
$1/\epsilon$, deep inside the corner.

\emph{Why a negative value contradicts SOS.} An exact identity
(\S\ref{sec:mech-identity}) writes the evaluated combination as a
part that is the same for every $Q$ minus a single $Q$-dependent term.
After scaling, the fixed part tends to a limiting tangent kernel whose
pairing with $c$ is below $-2^{42}$ by exact rational arithmetic
(Lemma~\ref{lem:witness}). Positivity of $Q$ forces the $Q$-dependent
term to be small, with constants independent of $Q$, of its size and of
its rank (Sections~\ref{sec:tails} and~\ref{sec:continuation}).
Theorem~\ref{thm:finite-scale} states the precise result: for every
integer $k\ge16$ and $m=2^{2k}$, if $\calB_m$ is a sum of squares, then
\[
 0\le\Phi_m<-2^{42}+2^{90}m^{-1/2}+2^{80}m^{-1/262144}.
\]
The right side is negative once $m$ is large enough. Then no positive
semidefinite Gram matrix, and hence no sum of squares, represents
$\calB_m$.

\begin{corollary}[Explicit threshold]\label{cor:threshold}
The corner $\calB_{m_0}$ is not a sum of squares for
$m_0=\CornerDepth$. Consequently $F_N$ is not a sum of squares for
every $N\ge2m_0=\ObstructionOrder$.
\end{corollary}
\begin{proof}
Set $k=38\cdot2^{17}+1=\ScaleExponent$ in the bound
\eqref{eq:finite-scale} of Theorem~\ref{thm:finite-scale}.
Convexity gives $2^{-x}\le1-x/2$ for $0\le x\le1$, so the right side is
\[
 -2^{42}+2^{90-k}+2^{42}2^{-2^{-17}}
 \le-2^{24}+2^{90-k}<0.
\]
This contradicts $\Phi_{m_0}\ge0$. Substitution preserves a sum of squares,
and the stabilized corner is a restriction of every $F_N$ with
$N\ge2m_0$.
\end{proof}

\subsection{Small orders and the open range}
For every $1\le N\le50$, $F_N$ is a sum of squares
(Proposition~\ref{prop:positive}). Here $F_1=0$; for $2\le N\le50$ the
statement follows from exact rational Gram certificates, and for
$2\le N\le20$ it is also formally proved in Lean~4 with Mathlib. The remaining orders are open:
\[
 51\le N<\ObstructionOrder.
\]
The smallest non-SOS order therefore lies between $51$ and
$\ObstructionOrder$, inclusive; monotonicity of the answer within this
interval is not known.
Section~\ref{sec:discussion} discusses an expectation for this range and the
evidence for it.

\subsection{Nonnegativity and sums of squares}\label{sec:context}
The B\"ottcher--Wenzel inequality
\[
 \|[X,Y]\|_F^2\le2\|X\|_F^2\|Y\|_F^2
\]
holds for arbitrary real square matrices
\cite{bottcher2005,laszlo2007,vong2008,bottcher2008}. If $X\ne0$,
replacing $Y$ by $Y_\perp=Y-\langle X,Y\rangle_F X/\|X\|_F^2$ does not change the
commutator. Expanding $\|Y_\perp\|_F^2$ in the resulting inequality gives
exactly $F_N\ge0$; the case $X=0$ is immediate. Thus the nonnegativity
assertion of Theorem~\ref{thm:negative} holds without the Toeplitz
hypothesis.

A nonnegative quartic need not be a sum of squares of quadratic forms.
This happens already for biquadratic forms, which have degree two in
each of two groups of variables, as $F_N$ has in $x$ and in $y$: Choi
gave a nonnegative biquadratic form in $3+3$ variables that is not a sum
of squares \cite{choi1975}. A form is a sum of squares exactly when it
admits a positive semidefinite Gram matrix, and its Gram matrices form an
affine family \cite{choi1995}; this Gram-matrix method underlies both the
certificates and the obstruction in this paper. L\'aszl\'o formulated
the Toeplitz question after studying the sum-of-squares problem for
the B\"ottcher--Wenzel form of general matrices
\cite{laszlo2010,laszlo2012};
Lu and Wenzel discuss it together with the broader commutator
inequalities \cite{lu2017}. The related DDVV inequalities involve sums
of commutator norms for families of symmetric matrices
\cite{desmet1999,ge2008,lu2011jfa}. Their sum-of-squares questions are
likewise distinct from the inequalities themselves; a recent survey
places both in this setting \cite{ge2024published}.

Theorem~\ref{thm:negative} shows that the Toeplitz structure does not
close the gap between the two properties. A sum of squares requires a
positive semidefinite Gram matrix on the whole space of wedge
coordinates, whereas nonnegativity only tests decomposable wedges
$x\wedge y$. The proof finds a negative direction for the stronger
requirement in a limiting kernel and then controls every error at a
finite order.

\subsection{Verification}
The negative conclusion of Theorem~\ref{thm:negative} is formally proved
in Lean~4 with Mathlib, for the literal polynomial $F_N$ and arbitrary
finite sums of real homogeneous quadratic squares, with no remaining
analytic hypothesis. The proof includes the complete corner extraction,
the analytic estimates, and a kernel-checked exact witness
(Lemma~\ref{lem:witness}). Its expanded statement is
\path|ToeplitzSOS.Negative.sharp_negative_explicit|; the negation of the
catalogue conjecture is \path|ToeplitzSOS.Negative.not_MI15|.
The positive range is formally proved through order $20$.
Appendix~\ref{app:verification} gives the declarations, the scope of each
formal component, and the relation to the written proof below.
